\documentclass[11pt]{article}
\usepackage[margin=1in]{geometry}
\usepackage{amsmath}
\usepackage{booktabs}
\usepackage{tabularx}
\usepackage{float}
\usepackage{enumitem}
\usepackage{microtype}
\usepackage{xurl}
\usepackage[hidelinks]{hyperref}
\sloppy

\setlist{itemsep=2pt, topsep=4pt}
\newcommand{\col}[1]{\texttt{#1}}
\newcolumntype{Y}{>{\raggedright\arraybackslash}X}

\title{Missouri Flood Impact Explorer: App Guide}
\author{CSE 472 -- Project 2}
\date{November 8, 2026}

\begin{document}
\maketitle

\section*{Purpose}

The Missouri Flood Impact Explorer is a Google Earth Engine app. It shows the data behind the flood
impact model and the model's predictions, for each of Missouri's 115 counties and each month from
2017 to 2023. The model predicts monthly flood property loss so that counties can be ranked for
emergency resources. This guide summarizes the model inputs, then explains what the map shows for
each choice in the app.

%----------------------------------------------------------------
\section{Summary of the Reduced Inputs}

The model uses \textbf{12 inputs plus calendar month}. Each input stands for one way that flooding
causes loss. State is also used in the full Midwest model, but it drops out here because every
county is in Missouri.

\begin{table}[H]
\centering
\small
\begin{tabularx}{\textwidth}{llYl}
\toprule
\textbf{Block} & \textbf{Variable} & \textbf{Meaning} & \textbf{Source} \\
\midrule
Hazard, & \col{p\_month} & Total rain in the month & PRISM \\
same month & \col{rx5\_rl\_ratio} & Wettest 5 days $\div$ the county's 10-year 5-day level & PRISM \\
 & \col{snowmelt\_mm} & Snowmelt in the month & ERA5-Land \\
 & \col{frost\_days} & Days with minimum temperature below 0\,$^\circ$C & PRISM \\
\midrule
Hazard, & \col{sm\_root\_pct} & Soil wetness on day 1, as a percentile against 1950--1999 & ERA5-Land \\
day 1 & \col{swe\_mm} & Snowpack on day 1 of the month & ERA5-Land \\
\midrule
Rain climate & \col{c\_rx5\_rl10\_mm} & 10-year 5-day rain level, from 2000--2009 & PRISM \\
\midrule
Exposure and & \col{log\_built\_surface} & Built-up area, natural log, previous year & GHSL \\
susceptibility & \col{impervious\_mean} & Paved share of the county, previous year & Annual NLCD \\
 & \col{hand\_lt5\_share} & Share of land less than 5 m above nearest drainage & MERIT Hydro \\
 & \col{log\_land\_km2} & County land area, natural log & Census TIGER \\
\midrule
Adaptation & \col{wetland\_share} & Wetland share of the county, previous year & Annual NLCD \\
\bottomrule
\end{tabularx}
\end{table}

The model has two parts: one predicts the chance of any damage, and the other predicts how large
the loss would be if damage occurs. \textbf{Expected loss} is the two multiplied together, and
counties are ranked by that number.

%----------------------------------------------------------------
\section{How to Use the App}

\begin{table}[H]
\centering
\small
\begin{tabularx}{\textwidth}{lY}
\toprule
\textbf{Control} & \textbf{What it does} \\
\midrule
First Layer menu & Picks a group: predictions, county inputs, or a source dataset \\
Second Layer menu & Picks one layer within that group \\
Year and Month menus & Pick the month, from January 2017 to December 2023 \\
Previous and Next & Move one month at a time \\
Outline checkbox & Draws a cyan outline around the 6 highest-risk counties that month (the top 5\%) \\
Colour bar & Shows the colour range and units of the current layer \\
Click on a county & Outlines it in black and opens its details in the side panel \\
\bottomrule
\end{tabularx}
\end{table}

The app opens on \textbf{Model predictions, Expected loss, April 2017}, Missouri's largest damage
month in the test years.

\subsection*{Two kinds of map}

\begin{itemize}
  \item \textbf{County maps} give each county one colour. These are the values the model used,
        or the model's predictions.
  \item \textbf{Source maps} show the original gridded data inside Missouri, with county borders
        on top. Colour can vary within a county.
\end{itemize}

Values beyond the ends of the colour bar are drawn in the end colour.

%----------------------------------------------------------------
\section{What the Map Shows for Each Choice}

The last column of each table shows what makes the map change. ``Fixed'' means the map stays the
same when the date changes.

\subsection{Model predictions}

\begin{table}[H]
\centering
\small
\begin{tabularx}{\textwidth}{lYll}
\toprule
\textbf{Layer} & \textbf{What the map shows} & \textbf{Colour range} & \textbf{Changes by} \\
\midrule
Expected loss & Predicted dollars of loss in that county and month. Darker red is a higher predicted loss. This is the main heatmap. & \$100 to \$10M, log scale & Month \\
Chance of any damage & Probability that the county records any flood damage & 0\% to 60\% & Month \\
Rank within the state & The county's position among all 115. Darkest is rank 1, the highest expected loss. & 1 to 115 & Month \\
Recorded loss (SHELDUS) & The loss that was actually recorded, for comparison. Blank counties have no record that month. & \$1k to \$100M, log scale & Month \\
\bottomrule
\end{tabularx}
\end{table}

\subsection{Model inputs (county values)}

These are the exact values given to the model for that county and month.

\begin{table}[H]
\centering
\small
\begin{tabularx}{\textwidth}{>{\raggedright\arraybackslash}p{3.1cm}Yll}
\toprule
\textbf{Layer} & \textbf{What the map shows} & \textbf{Colour range} & \textbf{Changes by} \\
\midrule
Monthly rain & \col{p\_month}: county-average rain total & 0 to 300 mm & Month \\
Wettest 5 days / 10-year level & \col{rx5\_rl\_ratio}: how extreme the month's heaviest 5 days were. A value of 1 equals a 10-year event. & 0 to 1.3 & Month \\
Snowmelt & \col{snowmelt\_mm}: melt during the month & 0 to 35 mm & Month \\
Frost days & \col{frost\_days}: days with minimum temperature below 0\,$^\circ$C & 0 to 30 days & Month \\
Soil wetness percentile & \col{sm\_root\_pct}: how wet the soil was on day 1, compared with the same month in 1950--1999 & 0\% to 100\% & Month \\
Snowpack & \col{swe\_mm}: water held in snow on day 1 & 0 to 15 mm & Month \\
10-year 5-day rain & \col{c\_rx5\_rl10\_mm}: the county's long-run heavy-rain level & 125 to 210 mm & Fixed \\
Built-up surface (log) & \col{log\_built\_surface}: amount of built-up area, previous year & 14.5 to 18.6 & Year \\
Paved share & \col{impervious\_mean}: paved percentage of the county, previous year & 0\% to 25\% & Year \\
Land under 5 m above drainage & \col{hand\_lt5\_share}: share of the county that sits low near water & 10\% to 100\% & Fixed \\
County land area (log) & \col{log\_land\_km2}: size of the county & 6.5 to 7.9 & Fixed \\
Wetland share & \col{wetland\_share}: wetland share of the county, previous year & 0\% to 8.5\% & Year \\
\bottomrule
\end{tabularx}
\end{table}

\subsection{Source datasets}

\begin{table}[H]
\centering
\small
\begin{tabularx}{\textwidth}{l>{\raggedright\arraybackslash}p{2.8cm}Yll}
\toprule
\textbf{Group} & \textbf{Layer} & \textbf{What the map shows} & \textbf{Colour range} & \textbf{Changes by} \\
\midrule
PRISM & Monthly rain & Sum of daily rain & 0 to 300 mm & Month \\
 & Wettest single day & Largest one-day rain in the month & 0 to 100 mm & Month \\
 & Frost days & Days with minimum temperature below 0\,$^\circ$C & 0 to 31 days & Month \\
\midrule
ERA5-Land & Root-zone soil water & Soil water from 0 to 100 cm on day 1 & 0.15 to 0.45 m$^3$/m$^3$ & Month \\
 & Snowpack & Snow water on day 1; blank where there is none & 0 to 30 mm & Month \\
 & Snowmelt & Monthly melt; blank where there is none & 0 to 40 mm & Month \\
\midrule
GHSL & Built-up surface & Built area in each 100 m cell; blank where there is none & 0 to 3,000 m$^2$ & Year* \\
\midrule
NLCD & Land cover & 16 land-cover classes, listed in the panel & Classes & Year \\
 & Impervious surface & Paved percentage; blank where there is none & 0\% to 80\% & Year \\
 & Wetlands & Wetland areas in blue & Yes only & Year \\
\midrule
Terrain & Elevation & Height of the land & 70 to 500 m & Fixed \\
 & Slope & Steepness of the land & 0 to 15$^\circ$ & Fixed \\
 & Height above drainage & Height above the nearest stream. Darker is closer to stream level. & 0 to 30 m & Fixed \\
 & Land under 5 m above drainage & Low land near streams in blue & Yes only & Fixed \\
\bottomrule
\end{tabularx}
\end{table}

\noindent NLCD and GHSL use the previous year's data, so that the inputs are known before the month
being predicted. *GHSL only has 5-year images: the map shows the 2015 image for 2017--2018 and the
2020 image for 2019--2023. The county input is interpolated between images, so it can differ.

\subsection*{How a county input relates to its source map}

\begin{table}[H]
\centering
\small
\begin{tabularx}{\textwidth}{YY}
\toprule
\textbf{Source map} & \textbf{County input in the model} \\
\midrule
PRISM Monthly rain & County average of the same data (\col{p\_month}) \\
PRISM Frost days & Days calculated from the county-average daily temperature (\col{frost\_days}) \\
ERA5-Land Root-zone soil water & Converted to a percentile against 1950--1999 (\col{sm\_root\_pct}) \\
ERA5-Land Snowpack, Snowmelt & County average (\col{swe\_mm}, \col{snowmelt\_mm}) \\
GHSL Built-up surface & County total, then natural log (\col{log\_built\_surface}) \\
NLCD Impervious surface & County average (\col{impervious\_mean}) \\
NLCD Wetlands & County share (\col{wetland\_share}) \\
Terrain: Land under 5 m above drainage & County share (\col{hand\_lt5\_share}) \\
\midrule
PRISM Wettest single day; NLCD Land cover; Terrain Elevation, Slope, Height above drainage &
Shown for context. These layers are not model inputs themselves. \\
\bottomrule
\end{tabularx}
\end{table}

%----------------------------------------------------------------
\section{What Clicking a County Shows}

\begin{itemize}
  \item \textbf{Prediction:} chance of any damage, expected loss, loss if damage occurs, rank in
        the state, and recorded loss.
  \item \textbf{Model inputs:} the 12 values for that month. Built-up surface and land area are
        converted from logs to km$^2$ so they are easier to read.
  \item \textbf{County mean:} for most source maps, an extra row gives that map's average over the
        county. Land cover and the two yes-or-no maps (Wetlands and Land under 5 m) do not have one.
  \item \textbf{Chart:} every month from 2017 to 2023. The red line is the predicted expected loss
        and purple dots are recorded losses. The axis is in $\log_{10}$ dollars, so 3 is \$1k and
        6 is \$1M.
\end{itemize}

The details update when the month or layer changes.

%----------------------------------------------------------------
\section{Worked Example: April 2017}

With \textbf{Model predictions, Expected loss, April 2017}, the darkest counties are in southern
Missouri. The 6 outlined counties are below, and each had recorded damage. Together they hold
63\% of the state's \$150M recorded loss that month. In all, 35 of the 115 counties recorded damage.

\begin{table}[H]
\centering
\small
\begin{tabular}{lcccc}
\toprule
\textbf{County} & \textbf{Rank} & \textbf{Chance of damage} & \textbf{Expected loss} & \textbf{Recorded loss} \\
\midrule
Barry & 1 & 82\% & \$10.8M & \$1.3M \\
Taney & 2 & 83\% & \$9.7M & \$12.8M \\
Greene & 3 & 76\% & \$9.6M & \$0.6M \\
Phelps & 4 & 90\% & \$8.4M & \$1.3M \\
Howell & 5 & 87\% & \$8.2M & \$76.9M \\
Newton & 6 & 78\% & \$8.2M & \$1.3M \\
\bottomrule
\end{tabular}
\end{table}

\noindent The model picked the right area, but the dollar amounts differ widely. Howell recorded
about 9 times its expected loss, and Greene about one-fifteenth of it.

\textbf{To see why,} keep April 2017 and try these combinations:

\begin{table}[H]
\centering
\small
\begin{tabularx}{\textwidth}{lY}
\toprule
\textbf{Choice} & \textbf{What to look for} \\
\midrule
Model predictions, Recorded loss & Compare recorded damage with the predicted pattern \\
Model inputs, Monthly rain & The heaviest county rain; Howell had 430 mm \\
Model inputs, Wettest 5 days / 10-year level & Values above 1 mean worse than a 10-year event; Howell reached 1.58 \\
PRISM, Monthly rain & Where within counties the rain fell \\
Terrain, Land under 5 m above drainage & Where low land near streams sits \\
GHSL, Built-up surface & Where buildings are concentrated \\
\bottomrule
\end{tabularx}
\end{table}

%----------------------------------------------------------------
\section{Limits to Keep in Mind}

\begin{itemize}
  \item \textbf{Predictions are computed outside Earth Engine.} LightGBM cannot run in Earth Engine,
        so the predictions are uploaded as a table and do not change with the map.
  \item \textbf{They are not true forecasts.} Predictions use the rain that actually fell that
        month.
  \item \textbf{Compare counties more than dollar amounts.} Dollar estimates are typically off by
        several times. Missouri's top 5\% is only 6 counties, so one large month can swing results.
  \item \textbf{Only 2017--2023 is shown.} The model was fitted on 2000--2016, so these are the
        years it did not train on.
  \item \textbf{Blank does not always mean zero loss.} A county with no SHELDUS record may still
        have had unreported damage.
  \item \textbf{NFIP flood insurance and buyouts are not included.} That data is still being collected.
\end{itemize}

%----------------------------------------------------------------
\section{How the Predictions Are Calculated}

The app does not calculate predictions itself. It looks up a row in \col{mo\_flood\_app\_table.csv},
which is built ahead of time by \col{build\_app\_table.py} using the LightGBM hurdle model in
\col{reduced\_benchmark.py}. This section follows one county through that calculation:
\textbf{Miller County, April 2017}.

\subsection*{The two models}

Both models are trained only on Missouri counties for 2000--2016, so April 2017 is out of sample.

\begin{table}[H]
\centering
\small
\begin{tabularx}{\textwidth}{lYY}
\toprule
 & \textbf{Part 1: classifier} & \textbf{Part 2: regressor} \\
\midrule
Predicts & Chance of any damage that month, $p$ & $\log_{10}$ of the loss, assuming damage happened, $\mu$ \\
Trained on & All county-months & Only county-months with damage \\
Miller County & $p = 0.375$ & Not saved in the table \\
\bottomrule
\end{tabularx}
\end{table}

\noindent The number of trees is set by early stopping (fit on 2000--2012, validate on 2013--2016),
then multiplied by 1.15 for the final fit on all of 2000--2016.

\subsection*{Inputs for Miller County, April 2017}

\begin{table}[H]
\centering
\small
\begin{tabular}{ll}
\toprule
\textbf{Input} & \textbf{Value} \\
\midrule
Monthly rain & 250.8 mm \\
Wettest 5 days / 10-year level & 1.24 (above a 10-year event) \\
Snowmelt, snowpack, frost days & About 0 \\
Soil wetness percentile & 72\% \\
10-year 5-day rain & 142.6 mm \\
Log built-up surface & 16.02 \\
Paved share & 1.2\% \\
Land under 5 m above drainage & 18\% \\
Log land area & 7.34 \\
Wetland share & 0.35\% \\
Calendar month & April \\
\bottomrule
\end{tabular}
\end{table}

\noindent Calendar month is a 13th input that the side panel's input list does not show.

\subsection*{How each number in the side panel is produced}

The expected loss combines both models:
\[
\text{Expected loss} = p \times 10^{\mu} \times \text{smear}
\]
The smear factor is a single correction: the average of $10^{(\text{actual} - \text{predicted})}$
over the training months with damage. Without it, converting from $\log_{10}$ back to dollars would
give something closer to a typical loss and understate the average.

\begin{table}[H]
\centering
\small
\begin{tabularx}{\textwidth}{>{\raggedright\arraybackslash}p{3.6cm}Yl}
\toprule
\textbf{Side panel field} & \textbf{Calculation} & \textbf{Miller County} \\
\midrule
Chance of any damage & $p$ from Part 1 & 0.3753, shown as 38\% \\
Expected loss & $p \times 10^{\mu} \times \text{smear}$ & \$1,377,640, shown as \$1.4M \\
Loss if damage occurs & Expected loss $\div\ p$, which equals $10^{\mu} \times \text{smear}$ & \$3,670,660, shown as \$3.7M \\
Rank in the state & Position by expected loss among the 115 counties that month & 40 \\
Recorded loss (SHELDUS) & Actual inflation-adjusted loss; not a prediction or an input & \$2,563,200, shown as \$2.6M \\
\bottomrule
\end{tabularx}
\end{table}

\subsection*{Reading these numbers}

\begin{itemize}
  \item \textbf{Compare \$3.7M to the recorded loss, not \$1.4M.} Damage did occur, so the loss if
        damage occurs is the like-for-like estimate. Here it is about 1.4 times too high. The
        \$1.4M is an average that includes the 62\% chance of no damage.
  \item \textbf{Loss if damage occurs is closer to a mean than a typical value.} The smear factor
        pulls it toward rare, very large losses, so a typical damaging month would be lower.
\end{itemize}

\end{document}
